% Mechanical relocation generated from the preservation ledger.
% Existing prose is included byte-for-byte from preserved blocks.

\section{Modeling Elements}

The following sections document the mathematical objects and modeling assumptions used internally by the software.


\subsection{Spatial Representation}

Let $\mathcal{S}$ denote the set of stops in the area of interest. Each stop can act as an origin, a destination, or a transfer location.

\begin{definition}[OD pair]
An origin--destination pair is a tuple
\[
q = (s_o, s_d) \in \mathcal{Q} = \mathcal{S} \times \mathcal{S}.
\]
\end{definition}


\subsection{Temporal Representation}

The analysis period is partitioned into a set $\mathcal{T}$ of departure time intervals. Each interval represents a desired departure time window.

\begin{definition}[Departure time interval]
Each time bin $t\in\mathcal{T}$ corresponds to a half-open interval
\[
[a_t,b_t)\subset\mathbb{R},
\qquad a_t<b_t,
\]
where $a_t$ and $b_t$ denote the start and end time (in minutes) of the
desired departure interval. The end of the final bin may be included by an
explicit analysis-period boundary convention. Half-open membership prevents a
departure on a shared boundary from belonging to two adjacent bins.
\end{definition}

For each departure interval $t\in\mathcal{T}$, the desired departure window is the
interval $[a_t,b_t)$. The admissible access window introduced below is a
separate, closed eligibility interval and does not change this unique bin
membership.

\begin{definition}[OD demand]
For each OD pair $q \in \mathcal{Q}$ and departure time interval $t \in \mathcal{T}$, let
\[
f_q^t \ge 0
\]
denote the number of passengers wishing to depart during interval $t$ from origin $s_o$ to destination $s_d$.
\end{definition}

The full demand vector is denoted by
\[
\bm{f} = \{ f_q^t : q \in \mathcal{Q},\; t \in \mathcal{T} \}.
\]



\subsection{Passenger Journeys, Vehicle Legs, and Events}

OD demand represents passenger journeys from a first boarding at the origin to
a final alighting at the destination. A journey with $L$ vehicle legs contains
$L-1$ transfers and generates one boarding and one alighting on every leg. It
therefore contributes one initial boarding, one final alighting, $L-1$ transfer
boardings, and $L-1$ transfer alightings. Raw boarding and alighting totals are
not journey-origin and journey-destination marginals when transfers occur.

For a conditional destination model, $Q_{ot}$ denotes journeys whose first
boarding occurs at origin $o$ in departure period $t$. It must come from an
external journey-entry source or a declared low-dimensional production model;
it cannot silently be equated to all boardings at $(o,t)$. Every model must
state which productions, attractiveness values, routing shares, detection
factors, and dispersion parameters are fixed, normalized, estimated, or
regularized.

\subsection{Physical Stops and Timetable Identity}

Scenario-stop identifiers define the external OD and measurement geography. An
explicit mapping may aggregate platforms into physical stop places for transfer
construction; without one, the mapping is the identity. Such a mapping must
cover every scenario stop exactly once and is part of the model assumptions.
A route pattern is identified by line, direction, and its complete ordered
physical-stop sequence. Scheduled trips retain their exact event times, so
opposite directions, express patterns, platform variants, and different service
periods remain distinct unless an explicit mapping states otherwise.

\subsection{Observation Resolution}

The atomic observation is a boarding or alighting count attached to an
unambiguous scheduled event. Stop-, route-, period-, or vehicle-level totals
are derived aggregates with an explicit mapping. An absent record is
unobserved, whereas a recorded zero is an observed zero. This distinction is
preserved in the likelihood.

\subsection{Free, Frozen, and Structurally Zero Demand}

Let $\mathcal U$ denote the free OD-time cells and $\mathcal F$ the frozen
cells. Only $\bm f_{\mathcal U}$ belongs to the estimation parameterization;
fixed values $\bm f_{\mathcal F}$ are constants. Frozen-zero cells contribute
neither a parameter nor a forward-response column. Positive frozen cells
contribute a fixed observation offset. The complete demand vector is
reconstructed only when a complete reported OD table is required.

Structural zeros are frozen zeros justified before estimation by declared
timetable rules. A cell is excluded when no admissible path exists or when
every path exceeds the accepted transfer bound; optional rules may additionally
limit initial wait, journey time, or service frequency. Classification is
performed on OD-time candidates using the same access, transfer, and physical-
stop assumptions as routing. A nonzero pre-existing fixed value that conflicts
with a detected structural zero is an error rather than an implicit overwrite.

Connectivity is therefore defined at the spatial resolution of the scenario
stops used by the OD table. These identifiers may be individual platforms. A
pair that is disconnected under the identity mapping may become connected when
an authoritative or reviewed platform-to-physical-stop mapping supplies the
interchange. Even at physical-stop level, connectivity remains time dependent:
directionality, service hours, access windows, minimum transfer time, and the
transfer bound may make a particular OD-time cell infeasible. ``No path'' and
``a path exists but requires too many transfers'' are recorded as distinct
structural-zero reasons.
